We want to find out which Kolumbus lines, stops and hours are least punctual. A planner at
Kolumbus or in the county could use this to decide which timetables to fix first. We picked
this data because we use these buses ourselves, and because it forced real decisions: the API
only shows the current state, and every answer repeats all running trips.

Some things are not part of the project:
\begin{itemize}
  \item \textbf{Live (streaming) results.} A report once a night is enough for this question.
  \item \textbf{Trips that are not in the live feed at all.} We can only measure trips that
        Entur sends. Finding missing trips would need the full timetable (GTFS).
  \item \textbf{Passenger numbers and a finished dashboard.} Gold is read with Spark SQL and
        charts (Section~\ref{sec:results}).
\end{itemize}

We wrote the layers as Python modules with one entry point (\code{jobs/run\_pipeline.py})
instead of notebooks, because modules are easier to schedule and to test.
